\documentclass[11pt,a4paper]{article}
\usepackage[utf8]{inputenc}
\usepackage[T1]{fontenc}
\usepackage[french]{babel}
\usepackage{amsmath}
\usepackage{amssymb}
\usepackage{booktabs}
\usepackage{enumitem}
\usepackage[a4paper,margin=2.2cm]{geometry}

\title{Le logarithme décimal\\[2pt]\large Découverte : d'où il vient, à quoi il sert}
\author{Terminale, spécialité mathématiques}
\date{Version imprimable générée le 03/11/2026\\[2pt]\small à compiler avec pdflatex (TeXStudio : touche F5)}

\begin{document}
\maketitle

\noindent Ce document est la version imprimable du chapitre interactif \emph{Découverte}. Les ateliers manipulables (curseurs, table de logarithmes, quiz) restent dans la version numérique : ici, vous trouverez le contenu du chapitre, prêt à être annoté.

\section{Il y a 400 ans : un problème de calcul}
Au début du XVII\textsuperscript{e} siècle, les astronomes, les navigateurs et les marchands doivent multiplier et diviser des nombres énormes, à la main, avec beaucoup de décimales. Une seule multiplication de deux nombres de dix chiffres peut prendre plusieurs minutes, et une erreur se glisse vite : les astronomes passent l'essentiel de leur temps à calculer, pas à observer le ciel.

\paragraph{Repères historiques}
\begin{itemize}[leftmargin=1.5em,itemsep=2pt]
\item \textbf{1614} --- L'Écossais John Napier publie \emph{Mirifici logarithmorum canonis descriptio} : il invente les logarithmes précisément pour alléger ces calculs.
\item \textbf{1617} --- L'Anglais Henry Briggs propose d'utiliser la base 10 : ce sont les \textbf{logarithmes décimaux}, ceux de votre calculatrice.
\item \textbf{1624} --- Briggs publie ses tables : les logarithmes des nombres de 1 à 20\,000 et de 90\,000 à 100\,000, avec 14 décimales.
\item \textbf{1620--1630} --- Apparition de la \textbf{règle à calcul} (Gunter, puis Oughtred) : une échelle logarithmique que l'on fait glisser.
\item \textbf{XVIII\textsuperscript{e} siècle} --- \textbf{Euler} relie logarithme et exponentielle : deux fonctions réciproques. C'est la vision moderne.
\item \textbf{1972} --- La calculatrice scientifique (HP-35) rend les tables de logarithmes inutiles, après trois siècles et demi de service.
\end{itemize}

On prête à Laplace cette formule : «\,Par la réduction du travail, le logarithme double la vie de l'astronome.\,» Étymologie : \emph{logarithme} vient du grec \emph{logos} (rapport, raison) et \emph{arithmos} (nombre).

\section{L'idée géniale : additionner les exposants}
Napier part d'une propriété des puissances :
\[ 2^{p} \times 2^{q} = 2^{p+q}. \]
Pour \textbf{multiplier} deux puissances d'un même nombre, il suffit d'\textbf{additionner} leurs exposants : une multiplication devient une addition. Avec la base 10 :
\[ 10^{3} \times 10^{5} = 10^{8}, \qquad 1\,000 \times 100\,000 = 100\,000\,000. \]
Le nombre auquel on «\,relit\,» le résultat s'appelle le \textbf{logarithme} : pour 1\,000, c'est 3 ; pour 100\,000, c'est 5 ; pour leur produit, c'est 8. Autrement dit : \textbf{le logarithme transforme une multiplication en addition}.

\section{Première approche : log(b) est un exposant}
\textbf{log(b) est l'exposant qu'il faut mettre sur 10 pour obtenir b} : log(b) répond à la question «\,10 puissance combien donne b ?\,»

\begin{center}
\begin{tabular}{lccccccc}
\toprule
$b$ & 0,01 & 0,1 & 1 & 10 & 100 & 1\,000 & $10^{6}$ \\
$\log b$ & $-2$ & $-1$ & 0 & 1 & 2 & 3 & 6 \\
\bottomrule
\end{tabular}
\end{center}

Deux valeurs à ne jamais oublier : $\log 1 = 0$ et $\log 10 = 1$. Les nombres plus petits que 1 ont un logarithme \textbf{négatif}.

\paragraph{Lire la partie entière sans calculatrice.}
Si $10^{k} \leqslant b < 10^{k+1}$, alors $k \leqslant \log b < k+1$ : la partie entière de $\log b$ se lit en comptant les chiffres de $b$ (moins un). Exemple : $10^{2} \leqslant 573 < 10^{3}$, donc $2 \leqslant \log 573 < 3$ --- et c'est un encadrement, pas la valeur exacte.

\paragraph{Combien de chiffres a $2^{n}$ ?}
Comme $\log(2^{n}) = n \log 2$ et $\log 2 \approx 0{,}301$ :
\[ \log(2^{100}) = 100 \times 0{,}301 = 30{,}1, \]
donc $2^{100}$ est compris entre $10^{30}$ et $10^{31}$ : il possède \textbf{31 chiffres}. C'est exactement le type de question traitée au chapitre 4 du cours.

\section{Le logarithme aujourd'hui}
\begin{itemize}[leftmargin=1.5em,itemsep=2pt]
\item \textbf{Chimie (pH)} : un écart de 1 unité de pH correspond à une concentration 10 fois plus grande. Sans log, il faudrait écrire des nombres comme $0{,}0000000001$ mol/L.
\item \textbf{Acoustique (décibels)} : +10 dB = puissance sonore $\times 10$. Entre un chuchotement (30 dB) et un concert (90 dB), la puissance est multipliée par 1\,000\,000.
\item \textbf{Sismologie (Richter)} : +1 de magnitude $\approx$ amplitude des ondes $\times 10$ (et énergie $\times 32$ environ).
\item \textbf{Astronomie} : de la taille d'un atome ($10^{-10}$ m) à l'univers observable ($10^{26}$ m), il y a 36 ordres de grandeur --- une échelle logarithmique les fait tenir sur une seule page.
\item \textbf{Épidémies} : une population qui double à chaque étape devient une \textbf{droite} en échelle logarithmique : on lit la vitesse de propagation immédiatement.
\item \textbf{Informatique} : chercher un mot dans un dictionnaire d'un million de pages demande au plus 20 questions «\,avant/après\,», car $2^{20} > 10^{6}$ : les algorithmes de recherche sont dits «\,en $\log n$\,».
\item \textbf{Finance} : un placement qui rapporte 5\,\% par an double au bout d'environ 14 ans : les intérêts composés sont exponentiels.
\item \textbf{Ordres de grandeur} : «\,3 ordres de grandeur de plus\,» signifie $\times 1\,000$. Les physiciens, les ingénieurs et les économistes raisonnent en logarithmes.
\end{itemize}

\section{Trois idées à emporter}
\begin{enumerate}[leftmargin=1.6em,itemsep=2pt]
\item \textbf{log(b) est un exposant} : celui de 10 dans l'écriture de b. Rien de plus mystérieux.
\item \textbf{log transforme les multiplications en additions} : c'est ce qui a fait sa fortune pendant 350 ans.
\item \textbf{log compresse les échelles} : d'où sa présence dans le pH, les décibels, la magnitude des séismes, les graphiques d'épidémies.
\end{enumerate}

\newpage
\section{Mon carnet de bord}
Réponses personnelles consignées dans la version interactive --- sessions de travail : 03/11/2026.

\subsection*{Historique --- Pourquoi le logarithme a-t-il été inventé au XVII\textsuperscript{e} siècle ? Par qui, et pour résoudre quel problème concret ?}
\emph{Pas encore de réponse.}

\subsection*{Vie quotidienne --- Citez au moins trois situations de la vie quotidienne où le logarithme intervient (directement ou à travers une unité de mesure). Expliquez brièvement.}
\emph{Pas encore de réponse.}

\subsection*{Monde industriel --- Comment le logarithme aide-t-il les ingénieurs et les techniciens ? (échelles, tolérances, ordres de grandeur, graphiques\ldots{})}
\emph{Pas encore de réponse.}

\subsection*{Monde scientifique --- Pourquoi les scientifiques utilisent-ils des échelles logarithmiques ? Donnez des exemples (pH, décibels, Richter, astronomie\ldots{})}
\emph{Pas encore de réponse.}

\subsection*{Informatique --- En quoi le logarithme est-il lié à la complexité des algorithmes et à la recherche d'information ?}
\emph{Pas encore de réponse.}

\subsection*{Réflexion personnelle --- Qu'avez-vous appris en découvrant ce document ? Qu'est-ce qui vous a surpris ?}
\emph{Pas encore de réponse.}
\end{document}